% GENERATED FILE. Edit canonical lessons, then run npm run generate:books.
% canonical-source-hash: fnv1a64:56f0dd260ad343a2
% canonical-lessons: GU-C514-dhilu, GU-C514-skru, GU-C514-kasvu, GU-C514-khiti, GU-C514-latkavvu

\chapter{Something Has Come Loose}
\label{ch:gu-something-has-come-loose}
\hlchaptermodality{514}

\begin{chapteropening}
\textbf{By the end of this chapter:} \emph{I can say what is loose and fix it.}

Complete the last lesson of chapter 514: I can say what is loose and fix it.
\end{chapteropening}

\section[\texorpdfstring{\gu{ઢીલું} (ḍhīlũ)}{ઢીલું (ḍhīlũ)}]{\gu{ઢીલું} (ḍhīlũ) — loose}

\label{lesson:GU-C514-dhilu}



\begin{quote}
Before the new one: say the Gujarati for brass, then the Gujarati for scrap.
\end{quote}

\subsection*{You'll want to know: \gu{ઢીલું}}
\textbf{\gu{ઢીલું}} — \emph{ḍhīlũ} — \textquotedblleft{}loose\textquotedblright{}.

\begin{quote}
\gu{આ} \gu{ખુરશીનો} \gu{એક} \gu{પગ} \gu{ઢીલો} \gu{છે}.
\end{quote}

\begin{quote}
\emph{ā khurśīno ek pag ḍhīlo che.}
\end{quote}

\begin{quote}
One leg of this chair is loose.
\end{quote}

\begin{cousinweb}
\textbf{\gu{ઢીલું}} (\emph{ḍhīlũ}) comes from \emph{*ḍhila}, \textquotedblleft{}slack\textquotedblright{}, a form CDIAL reconstructs, and CDIAL lists the Gujarati word, \textquotedblleft{}loose, slack\textquotedblright{} (CDIAL 5590). It agrees with its noun: \textbf{\gu{પગ} \gu{ઢીલો}}, \textbf{\gu{ગાંઠ} \gu{ઢીલી}}.
\end{cousinweb}

\begin{grammarlens}[title={What you've built}]
One more word for describing what you lived and felt.
\end{grammarlens}

\subsection*{Your turn}
\begin{itemize}
\raggedright
  \item \emph{Say it:} \emph{ḍhīlũ}
  \item \emph{Say it:} \emph{ḍhīlũ}, once more
  \item \emph{Say it:} \emph{ā khurśīno ek pag ḍhīlo che.} — One leg of this chair is loose.
  \item \emph{Recall:} say the Gujarati for brass, then the Gujarati for scrap, then say \emph{ḍhīlũ} again
\end{itemize}

\begin{tcolorbox}[breakable,colback=teal!4,colframe=teal!35!black,title={Before you move on}]
What does \gu{ઢીલું} mean? (\textquotedblleft{}Loose\textquotedblright{}.) Say it once more.
\end{tcolorbox}
\section[\texorpdfstring{\gu{સ્ક્રૂ} (skrū)}{સ્ક્રૂ (skrū)}]{\gu{સ્ક્રૂ} (skrū) — a screw}

\label{lesson:GU-C514-skru}



\begin{quote}
Before the new one: say the Gujarati for scrap, then the Gujarati for loose.
\end{quote}

\subsection*{You'll want to know: \gu{સ્ક્રૂ}}
\textbf{\gu{સ્ક્રૂ}} — \emph{skrū} — \textquotedblleft{}a screw\textquotedblright{}.

\begin{quote}
\gu{બારીનો} \gu{એક} \gu{સ્ક્રૂ} \gu{ખોવાઈ} \gu{ગયો}.
\end{quote}

\begin{quote}
\emph{bārīno ek skrū khovāī gayo.}
\end{quote}

\begin{quote}
One of the window's screws has gone missing.
\end{quote}

A masculine noun.

\begin{cousinweb}
\textbf{\gu{સ્ક્રૂ}} (\emph{skrū}) is borrowed from English \emph{screw}.
\end{cousinweb}

\begin{grammarlens}[title={What you've built}]
One more word for describing what you lived and felt.
\end{grammarlens}

\subsection*{Your turn}
\begin{itemize}
\raggedright
  \item \emph{Say it:} \emph{skrū}
  \item \emph{Say it:} \emph{skrū}, once more
  \item \emph{Say it:} \emph{bārīno ek skrū khovāī gayo.} — One of the window's screws has gone missing.
  \item \emph{Recall:} say the Gujarati for scrap, then the Gujarati for loose, then say \emph{skrū} again
\end{itemize}

\begin{tcolorbox}[breakable,colback=teal!4,colframe=teal!35!black,title={Before you move on}]
What does \gu{સ્ક્રૂ} mean? (\textquotedblleft{}A screw\textquotedblright{}.) Say it once more.
\end{tcolorbox}
\section[\texorpdfstring{\gu{કસવું} (kasvũ)}{કસવું (kasvũ)}]{\gu{કસવું} (kasvũ) — to tighten}

\label{lesson:GU-C514-kasvu}



\begin{quote}
Before the new one: say the Gujarati for loose, then the Gujarati for a screw.
\end{quote}

\subsection*{You'll want to know: \gu{કસવું}}
\textbf{\gu{કસવું}} — \emph{kasvũ} — \textquotedblleft{}to tighten\textquotedblright{}.

\begin{quote}
\gu{અરુણે} \gu{ખુરશીના} \gu{બધા} \gu{સ્ક્રૂ} \gu{કસ્યા}.
\end{quote}

\begin{quote}
\emph{Aruṇe khurśīnā badhā skrū kasyā.}
\end{quote}

\begin{quote}
Arun tightened all the screws of the chair.
\end{quote}

Past, with \gu{તેણે} (he or she …): \gu{તેણે} \gu{સ્ક્રૂ} \gu{કસ્યો}, \emph{teṇe skrū kasyo}.

\begin{cousinweb}
\textbf{\gu{કસવું}} (\emph{kasvũ}) is inherited from Sanskrit \emph{kárṣati}, \textquotedblleft{}draws, pulls\textquotedblright{}, and CDIAL lists the Gujarati verb, \textquotedblleft{}to tighten\textquotedblright{} (CDIAL 2908).
\end{cousinweb}

\begin{grammarlens}[title={What you've built}]
One more word for describing what you lived and felt.
\end{grammarlens}

\subsection*{Your turn}
\begin{itemize}
\raggedright
  \item \emph{Say it:} \emph{kasvũ}
  \item \emph{Say it:} \emph{kasvũ}, once more
  \item \emph{Say it:} \emph{Aruṇe khurśīnā badhā skrū kasyā.} — Arun tightened all the screws of the chair.
  \item \emph{Recall:} say the Gujarati for loose, then the Gujarati for a screw, then say \emph{kasvũ} again
\end{itemize}

\begin{tcolorbox}[breakable,colback=teal!4,colframe=teal!35!black,title={Before you move on}]
What does \gu{કસવું} mean? (\textquotedblleft{}To tighten\textquotedblright{}.) Say it once more.
\end{tcolorbox}
\section[\texorpdfstring{\gu{ખીંટી} (khī̃ṭī)}{ખીંટી (khī̃ṭī)}]{\gu{ખીંટી} (khī̃ṭī) — a wall peg}

\label{lesson:GU-C514-khiti}



\begin{quote}
Before the new one: say the Gujarati for a screw, then the Gujarati for to tighten.
\end{quote}

\subsection*{You'll want to know: \gu{ખીંટી}}
\textbf{\gu{ખીંટી}} — \emph{khī̃ṭī} — \textquotedblleft{}a wall peg\textquotedblright{}.

\begin{quote}
\gu{મારી} \gu{છત્રી} \gu{ખીંટી} \gu{પર} \gu{છે}.
\end{quote}

\begin{quote}
\emph{mārī chatrī khī̃ṭī par che.}
\end{quote}

\begin{quote}
My umbrella is on the peg.
\end{quote}

A feminine noun.

\begin{cousinweb}
\textbf{\gu{ખીંટી}} (\emph{khī̃ṭī}) is filed by CDIAL under Sanskrit \emph{kīla}, \textquotedblleft{}a stake, a peg\textquotedblright{}, with the Gujarati word glossed \textquotedblleft{}a stopper\textquotedblright{} (CDIAL 3202); on a wall it is a peg to hang things on.
\end{cousinweb}

\begin{grammarlens}[title={What you've built}]
One more word for describing what you lived and felt.
\end{grammarlens}

\subsection*{Your turn}
\begin{itemize}
\raggedright
  \item \emph{Say it:} \emph{khī̃ṭī}
  \item \emph{Say it:} \emph{khī̃ṭī}, once more
  \item \emph{Say it:} \emph{mārī chatrī khī̃ṭī par che.} — My umbrella is on the peg.
  \item \emph{Recall:} say the Gujarati for a screw, then the Gujarati for to tighten, then say \emph{khī̃ṭī} again
\end{itemize}

\begin{tcolorbox}[breakable,colback=teal!4,colframe=teal!35!black,title={Before you move on}]
What does \gu{ખીંટી} mean? (\textquotedblleft{}A wall peg\textquotedblright{}.) Say it once more.
\end{tcolorbox}
\section[\texorpdfstring{\gu{લટકાવવું} (laṭkāvvũ)}{લટકાવવું (laṭkāvvũ)}]{\gu{લટકાવવું} (laṭkāvvũ) — to hang (something) from}

\label{lesson:GU-C514-latkavvu}



\begin{quote}
Before the new one: say the Gujarati for to tighten, then the Gujarati for a wall peg.
\end{quote}

\subsection*{You'll want to know: \gu{લટકાવવું}}
\textbf{\gu{લટકાવવું}} — \emph{laṭkāvvũ} — \textquotedblleft{}to hang (something) from\textquotedblright{}.

\begin{quote}
\gu{મીરાએ} \gu{થેલી} \gu{ખીંટી} \gu{પર} \gu{લટકાવી}.
\end{quote}

\begin{quote}
\emph{Mīrāe thelī khī̃ṭī par laṭkāvī.}
\end{quote}

\begin{quote}
Mira hung the bag on the peg.
\end{quote}

Past, with \gu{તેણે} (he or she …): \gu{તેણે} \gu{થેલી} \gu{લટકાવી}, \emph{teṇe thelī laṭkāvī}.

\begin{cousinweb}
\textbf{\gu{લટકાવવું}} (\emph{laṭkāvvũ}) is built from \textbf{\gu{લટકવું}}, \textquotedblleft{}to hang, to dangle\textquotedblright{}, with the causative \emph{-āv-}: to make something hang. CDIAL files the family under \emph{*laṭṭa}, \textquotedblleft{}hanging\textquotedblright{} (CDIAL 10918).
\end{cousinweb}

\begin{grammarlens}[title={What you've built}]
One more word for describing what you lived and felt.
\end{grammarlens}

\subsection*{Your turn}
\begin{itemize}
\raggedright
  \item \emph{Say it:} \emph{laṭkāvvũ}
  \item \emph{Say it:} \emph{laṭkāvvũ}, once more
  \item \emph{Say it:} \emph{Mīrāe thelī khī̃ṭī par laṭkāvī.} — Mira hung the bag on the peg.
  \item \emph{Recall:} say the Gujarati for to tighten, then the Gujarati for a wall peg, then say \emph{laṭkāvvũ} again
\end{itemize}

\begin{tcolorbox}[breakable,colback=teal!4,colframe=teal!35!black,title={Before you move on}]
What does \gu{લટકાવવું} mean? (\textquotedblleft{}To hang (something) from\textquotedblright{}.) Say it once more.
\end{tcolorbox}
